\chapter{Kiểm thử và Triển khai (Testing \& Deployment)}

\section{Chiến lược Kiểm thử Frontend (Frontend Testing Strategy)}
Trong phát triển phần mềm hiện đại, việc đảm bảo chất lượng mã nguồn (Quality Assurance) không chỉ phụ thuộc vào việc kiểm tra thủ công (Manual Testing) mà cần được tự động hóa để phát hiện lỗi sớm nhất có thể. Đối với một ứng dụng Frontend có tính tương tác cao như Trello Mini, nhóm đã xây dựng một chiến lược kiểm thử bao gồm hai cấp độ chính: Kiểm thử Đơn vị (Unit Testing) và Kiểm thử Giao diện (Component Testing).

\subsection{Kiểm thử Đơn vị (Unit Testing) với Vitest}
Kiểm thử đơn vị tập trung vào việc xác minh tính đúng đắn của các hàm logic, các hooks tự tạo (Custom Hooks) và đặc biệt là logic quản lý trạng thái trong Zustand Store. Nhóm quyết định sử dụng \textbf{Vitest} thay vì Jest truyền thống, vì Vitest tương thích hoàn hảo với hệ sinh thái Vite, mang lại tốc độ thực thi cực nhanh và hỗ trợ sẵn TypeScript mà không cần cấu hình phức tạp.

Một ví dụ tiêu biểu là việc kiểm thử hàm xử lý logic khi người dùng kéo thả một thẻ (Card) sang một cột (List) mới. Hàm này yêu cầu tính toán lại chỉ mục (index) và cập nhật mảng dữ liệu mà không làm thay đổi các thẻ khác:

\begin{Verbatim}
import { describe, it, expect } from 'vitest';
import { useBoardStore } from '../store/boardStore';

describe('Board Store - Drag and Drop Logic', () => {
  it('should move card to a different list correctly', () => {
    // Khởi tạo trạng thái ban đầu
    const store = useBoardStore.getState();
    store.addList('List A');
    store.addList('List B');
    store.addCard('List A', 'Task 1');

    // Thực hiện hành động kéo thả
    store.moveCard(cardId, 'List A', 'List B', 0);

    // Kiểm tra kết quả
    const updatedState = useBoardStore.getState();
    expect(updatedState.lists.find(l => l.name === 'List A').cards).toHaveLength(0);
    expect(updatedState.lists.find(l => l.name === 'List B').cards[0].title).toBe('Task 1');
  });
});
\end{Verbatim}

\subsection{Kiểm thử Giao diện (Component Testing) với React Testing Library}
Trong khi Unit Test đảm bảo logic hoạt động đúng, \textbf{React Testing Library (RTL)} được sử dụng để kiểm tra xem giao diện có phản hồi chính xác với các tương tác của người dùng hay không. RTL khuyến khích cách viết test dựa trên những gì người dùng thực sự nhìn thấy và tương tác (như tìm kiếm một nút bấm qua văn bản hiển thị) thay vì dựa vào chi tiết triển khai bên trong (như tên class CSS).

\begin{Verbatim}
import { render, screen, fireEvent } from '@testing-library/react';
import AddCardForm from './AddCardForm';

test('renders input field and adds new card on submit', () => {
  const mockAddCard = vi.fn();
  render(<AddCardForm listId="1" onAdd={mockAddCard} />);

  // Mô phỏng người dùng nhập text
  const input = screen.getByPlaceholderText('Nhập tiêu đề thẻ...');
  fireEvent.change(input, { target: { value: 'Công việc mới' } });

  // Mô phỏng click nút Thêm
  const button = screen.getByText('Thêm thẻ');
  fireEvent.click(button);

  // Xác minh hàm onAdd được gọi đúng tham số
  expect(mockAddCard).toHaveBeenCalledWith('1', 'Công việc mới');
});
\end{Verbatim}

\subsection{Kiểm thử đầu cuối (End-to-End Testing) với Playwright}
Trong khi Unit và Component Test kiểm tra các phần nhỏ của ứng dụng, Kiểm thử đầu cuối (E2E) mô phỏng toàn bộ hành trình của người dùng từ khi mở trình duyệt cho đến khi hoàn thành mục tiêu. Nhóm lựa chọn \textbf{Playwright} của Microsoft làm công cụ kiểm thử E2E chủ đạo.

\begin{itemize}
    \item \textbf{Lợi thế của Playwright:} Hỗ trợ kiểm thử trên nhiều trình duyệt khác nhau (Chromium, Firefox, WebKit) chỉ với một bộ mã duy nhất. Cơ chế \textbf{Auto-wait} giúp giảm thiểu tình trạng test bị "flickering" (lúc chạy được lúc không) do trễ mạng hoặc animation.
    \item \textbf{Kiểm thử kéo thả (Drag and Drop):} Một trong những thách thức lớn nhất của Trello Mini là kiểm thử tính năng kéo thả. Playwright cung cấp API \texttt{page.dragAndDrop(source, target)} giúp tự động hóa hoàn toàn thao tác này.
\end{itemize}

\begin{Verbatim}
// Ví dụ test case kéo thả một card sang cột khác bằng Playwright
test('should drag a card from Todo to In Progress list', async ({ page }) => {
  await page.goto('/board/1');

  const sourceCard = page.locator('.card-item:has-text("Lập kế hoạch")');
  const targetList = page.locator('.list-container:has-text("In Progress")');

  // Thực hiện hành động kéo thả
  await sourceCard.dragTo(targetList);

  // Kiểm chứng card đã nằm trong cột mới
  await expect(targetList).toContainText('Lập kế hoạch');
});
\end{Verbatim}

Việc áp dụng E2E Test giúp nhóm tự tin hơn mỗi khi thay đổi kiến trúc lớn mà không lo sợ làm hỏng các luồng nghiệp vụ cơ bản của người dùng.

\section{Các biện pháp bảo mật (Security Best Practices)}
An toàn thông tin là một khía cạnh không thể tách rời trong phát triển phần mềm, ngay cả với các ứng dụng Frontend.

\begin{itemize}
    \item \textbf{Ngăn chặn XSS (Cross-Site Scripting):} React mặc định sẽ tự động "escape" toàn bộ giá trị render ra giao diện để ngăn chặn việc thực thi mã Javascript độc hại từ phía người dùng. Nhóm tuyệt đối không sử dụng thuộc tính \texttt{dangerouslySetInnerHTML} trừ những trường hợp thực sự cần thiết và đã được kiểm duyệt kỹ càng.
    \item \textbf{Làm sạch dữ liệu đầu vào (Input Sanitization):} Đối với các nội dung do người dùng nhập vào (như mô tả thẻ công việc), nhóm sử dụng thư viện \textbf{DOMPurify} để loại bỏ các thẻ HTML nguy hiểm trước khi lưu trữ hoặc hiển thị.
    \item \textbf{Bảo vệ LocalStorage:} Nhóm nhận thức được rằng \texttt{LocalStorage} dễ bị tấn công bởi các đoạn mã độc. Vì vậy, ứng dụng không lưu trữ bất kỳ thông tin nhạy cảm nào như mật khẩu hay Token bí mật dưới dạng văn bản thuần túy. Đối với các hệ thống yêu cầu bảo mật cao, nhóm đề xuất sử dụng \textbf{HttpOnly Cookies} để lưu trữ phiên đăng nhập.
\end{itemize}

\section{Tích hợp và Triển khai liên tục (CI/CD)}

\subsection{Đóng gói ứng dụng (Build) với Vite}
Trước khi ứng dụng có thể được đưa lên môi trường sản xuất (Production), mã nguồn React (JSX/TypeScript) cần được biên dịch và đóng gói thành các tệp HTML, CSS và JavaScript thuần túy. Quá trình này được thực hiện thông qua \textbf{Vite} (với lõi là Rollup).

Vite thực hiện hàng loạt các tối ưu hóa trong quá trình build:
\begin{itemize}
    \item \textbf{Minification:} Nén và loại bỏ khoảng trắng, comments, các đoạn mã không được sử dụng (Tree-shaking) để giảm thiểu kích thước tệp.
    \item \textbf{Code Splitting:} Chia nhỏ bundle khổng lồ thành nhiều tệp tin (chunks) nhỏ hơn, cho phép trình duyệt tải bất đồng bộ (Lazy Load) khi cần thiết.
    \item \textbf{Cache Busting:} Tự động thêm mã băm (hash) vào tên tệp tin (ví dụ: \texttt{main.[hash].js}) để đảm bảo trình duyệt luôn tải phiên bản mới nhất thay vì sử dụng bộ nhớ đệm (cache) cũ.
\end{itemize}

\subsection{Tự động hóa triển khai bằng GitHub Actions và Vercel}
Để rút ngắn thời gian đưa sản phẩm tới tay người dùng và giảm thiểu sai sót do thao tác thủ công, nhóm đã thiết lập một luồng CI/CD (Continuous Integration / Continuous Deployment) hoàn toàn tự động sử dụng \textbf{GitHub Actions} kết hợp với nền tảng đám mây \textbf{Vercel}.

Quy trình hoạt động diễn ra theo các bước sau:
\begin{enumerate}
    \item \textbf{Push code:} Mỗi khi một thành viên trong nhóm đẩy (push) mã nguồn mới lên nhánh \texttt{main} trên GitHub, một workflow sẽ tự động được kích hoạt.
    \item \textbf{CI Pipeline (Kiểm tra chất lượng):} GitHub Actions sẽ khởi tạo một máy ảo (Runner), cài đặt các thư viện (\texttt{npm install}), chạy công cụ kiểm tra cú pháp (\texttt{ESLint}) và thực thi toàn bộ bộ Unit Test. Nếu bất kỳ một test case nào thất bại, quy trình sẽ dừng lại ngay lập tức và gửi thông báo cảnh báo cho nhóm.
    \item \textbf{CD Pipeline (Triển khai):} Nếu bước CI vượt qua thành công, mã nguồn sẽ được chuyển giao cho Vercel. Nền tảng này sẽ tự động chạy lệnh \texttt{npm run build} và phân phối (deploy) các tệp tĩnh (static files) lên mạng lưới máy chủ toàn cầu (CDN - Content Delivery Network).
\end{enumerate}

Nhờ mô hình này, ứng dụng luôn được đảm bảo ở trạng thái ổn định nhất mỗi khi phát hành. Người dùng có thể truy cập vào đường dẫn trực tiếp do Vercel cung cấp để trải nghiệm phiên bản mới nhất của Trello Mini chỉ trong vòng chưa đầy 2 phút sau khi mã nguồn được hợp nhất.

\section{Giám sát và Ghi vết (Monitoring \& Logging)}
Sau khi ứng dụng đã được triển khai, việc theo dõi hoạt động và phát hiện lỗi phát sinh từ phía người dùng là cực kỳ quan trọng.

\begin{itemize}
    \item \textbf{Sentry (Error Tracking):} Nhóm đề xuất tích hợp Sentry để tự động bắt và gửi thông báo về các lỗi Runtime xảy ra trên trình duyệt của người dùng. Mỗi khi có một lỗi xuất hiện, Sentry sẽ cung cấp đầy đủ thông tin về trình duyệt, hệ điều hành, dòng code bị lỗi và thậm chí là các bước người dùng đã thực hiện trước đó (Breadcrumbs).
    \item \textbf{LogRocket (Session Replay):} Để hiểu sâu hơn về trải nghiệm người dùng và tái hiện các lỗi khó (Edge cases), các công cụ như LogRocket cho phép "ghi hình" lại phiên làm việc của người dùng (dưới dạng các sự kiện DOM, không phải video thực tế để đảm bảo tính riêng tư). Điều này giúp các lập trình viên nhìn thấy chính xác những gì người dùng đã thấy khi gặp lỗi.
\end{itemize}
Việc kết hợp giám sát chủ động giúp duy trì chất lượng ứng dụng và cải thiện độ tin cậy của hệ thống theo thời gian.

